% Hoja 1: límites y continuidad.
%
% Una hoja de problemas tiene la misma forma que un tema; lo que cambia son
% las salidas. Del mismo fichero salen tres hojas distintas:
%
%   problems           los enunciados, con las pistas
%   problems-answers   y el resultado de cada problema
%   problems-teacher   y la solución entera y los criterios de corrección

\input{didacta-bootstrap}
\usepackage{didacta}

%% Solo un respaldo: Didacta inyecta estos datos desde `course.yaml` y
%% `year.yaml`, en el idioma que se compila. Ver tema-1.tex.
\DidactaCourse{
  title   = {Cálculo I},
  teacher = {Tu nombre},
}
\DidactaDocument{Hoja 1: límites y continuidad}

\begin{document}
%% Una hoja no quiere portada: quiere una cabecera en la primera página.
\DidactaSheetHeader

%% La composición, que el motor escribe a partir de `year.yaml`.

\DidactaProblem{calculo/limites/limites-de-cocientes}
\DidactaProblem{calculo/continuidad/discontinuidades-evitables}
\DidactaProblem{calculo/continuidad/tres-raices-por-bolzano}

\end{document}
